% Interface et tests utilisateurs (UAT), écrit à la main à partir des rapports v2/docs/lab/reports/ :
% 2026-11-14-prototype-decision.md, -prototype-carte.md, -uat-interface.md, -accueil-personas.md, -produit.md
\part{Interface et tests utilisateurs}

\section{Deux prototypes branchés sur le jumeau}
\begin{itemize}
  \item \textbf{V1, l'écran de décision} (\texttt{v2/prototype/}). C'est une page servie par un
        petit serveur Python qui appelle le jumeau ; aucun moteur n'est dupliqué en
        JavaScript. Elle montre :
        \begin{itemize}
          \item l'écart au scénario tendanciel, décomposé exactement par flux (entrées,
                diplômés étrangers, sorties, régulation) ;
          \item le délai d'effet de chaque levier ;
          \item une fourchette dessinée en hachures ;
          \item une matrice département × année qu'on peut retrier.
        \end{itemize}
        Neuf leviers du jumeau y sont rangés en quatre familles : \emph{former},
        \emph{attirer}, \emph{installer}, \emph{retenir}. Chacun porte deux badges, son délai
        d'effet et son niveau de preuve (sourcé ou supposé).
        \begin{itemize}
          \item Exemple vérifié : kinés, places en formation $+20$~\%, diplômés étrangers
                $\times 1{,}5$, cessation reportée d'un an et régulation vers le salariat. On
                obtient $+25\,007$ professionnels en 2050 par rapport à la tendance
                ($+15{,}6$~\%), dont entrées $+20\,665$, étrangers $+14\,765$, régulation
                $-10\,773$ et sorties $+339$ (\mesure).
          \item Recalcul : 77~ms aller-retour (\mesure).
        \end{itemize}
  \item \textbf{V2, la carte} (\texttt{v2/prototype-carte/}). C'est une carte MapLibre plein
        écran. Un mouvement de levier recolore les 101 départements et met à jour la fiche en
        53~ms en médiane, 66~ms au 95\textsuperscript{e} centile et 87~ms au maximum
        ($n=35$, \mesure), sous la cible de 150~ms.
\end{itemize}

\section{Méthode des tests utilisateurs}
Les tests ont été conduits par des \textbf{agents qui jouent des personas}, pas par de vrais
utilisateurs. C'est leur principale limite.
\begin{itemize}
  \item \textbf{Cinq personas} : une directrice générale d'ARS, l'administration centrale
        (DGOS), un·e élu·e, un utilisateur mobile et un testeur adversarial. Chacun avait une
        mission chronométrée.
  \item \textbf{Vérification de chaque constat} par un second agent chargé de le réfuter :
        46 agents pour la V1, 44 pour la V2.
  \item \textbf{Classement} des constats par gravité $\times$ fréquence.
  \item \textbf{Durées} : estimées par les agents, pas mesurées. Seules les latences en
        millisecondes sont mesurées.
\end{itemize}

\section{Résultats}
\paragraph{V1.}
\begin{itemize}
  \item 11 tâches réussies sur 22, et 6 échecs.
  \item 38 constats confirmés sur 40.
  \item Les échecs les plus graves :
        \begin{itemize}
          \item \textbf{lire l'effectif d'une année intermédiaire} (2035) : il n'existe pas de
                sélecteur d'horizon ;
          \item \textbf{descendre jusqu'à la commune} : la projection s'arrête au
                département ;
          \item \textbf{exporter une fiche} avec ses sources et ses hypothèses.
        \end{itemize}
\end{itemize}

\paragraph{V2.} Verdict « à porter en démo », mais la recette est \textbf{faussée} : les cinq
personas ont partagé un seul onglet de navigateur, faute de pouvoir en ouvrir d'autres. Au
moins 4 échecs sur 6 viennent de ce protocole, et non du produit. Les latences d'environ 1~s
observées ont été réfutées : elles venaient d'un onglet en arrière-plan.

\paragraph{Écran d'accueil.} Les trois personas principaux refusent l'accueil actuel (une
carte de France, les leviers ouverts, aucun chiffre avant un clic).
\begin{itemize}
  \item \textbf{Ce qu'ils veulent, tous les trois} :
        \begin{itemize}
          \item une phrase avec un chiffre dès l'ouverture (« kinés libéraux en ARA : 13\,240
                en 2026, environ 20\,100 en 2050 au tendanciel ») ;
          \item les leviers repliés derrière « Simuler une politique » ;
          \item la fourchette toujours visible ;
          \item le périmètre dit en clair.
        \end{itemize}
  \item \textbf{Ce qui les divise : le point d'entrée.}
        \begin{itemize}
          \item l'ARS veut entrer par sa région ;
          \item l'élu·e par son territoire, jusqu'à la commune ;
          \item l'Administration par la France entière, avec sources, millésime et version du
                modèle.
        \end{itemize}
  \item \textbf{Synthèse proposée} : un champ unique « Quel territoire ? » sur une carte
        estompée, et deux modes, « simple » (des phrases) et « expert » (flux, sources,
        export), sur un même modèle.
\end{itemize}

\paragraph{Retours directs de l'utilisateur·rice sur la carte.}
\begin{itemize}
  \item un zoom plus fort sur régions et départements ;
  \item des cases et des encarts plus grands ;
  \item des limites de régions ARS mieux séparées ;
  \item les départements et régions d'outre-mer en encarts ;
  \item un milieu de gamme de couleur qui ne se confond pas avec le fond.
\end{itemize}
Ils ont été intégrés au fil de sept itérations.

\section{Direction artistique}
Trois directions ont été proposées : sémiologie graphique (Bertin), gamme de signalisation, et
topographie en isolignes. La direction « sémiologie graphique » a été prototypée :
\begin{itemize}
  \item fond blanc ;
  \item une seule gamme de couleur par variable, jamais rouge ni verte (le lexique neutre rendu
        visible) ;
  \item l'incertitude en hachures ;
  \item une matrice qu'on peut réordonner.
\end{itemize}
Pour que les designers de l'équipe puissent reprendre la direction artistique dans Figma, une
source unique de \emph{tokens} au format W3C est proposée, où les directions deviendraient des
\emph{modes} d'une même collection. Elle n'est pas encore construite.

\section{Produit}
L'outil n'est \textbf{pas achetable en l'état} : sur 8 conditions de signature recensées,
aucune n'est remplie complètement. La voie réaliste est celle d'un commun numérique libre,
porté par une institution (DREES, une ARS pilote, beta.gouv), puis de prestations de moins de
60~k€ HT par ARS (seuil de dispense de procédure au 1\textsuperscript{er} avril 2026, cité
dans le rapport produit).

\section{Menaces à la validité}
\begin{itemize}
  \item Les personas sont des agents, pas des utilisateurs réels ; les durées sont estimées.
  \item La recette V2 a été faussée par l'onglet partagé, et aucune capture n'a été écrite sur
        disque.
  \item Les prototypes ont changé pendant les tests.
  \item Les chiffres de la V1 antérieurs aux corrections du jumeau ne sont plus valables (par
        exemple, quotas $+20$~\% chez les kinés : $+19\,110$ et non $+8\,645$), et le rapport
        le signale.
\end{itemize}
